\chapter{Kiến trúc và mô hình dữ liệu}
\label{ch:kien-truc-va-du-lieu}

\section{Các thành phần trong hệ thống}

\begin{table}[htbp]
\centering
\caption{Năm thành phần của hệ thống}
\label{tab:nam-service}
\begin{tabular}{@{}lL{4.9cm}L{2.4cm}L{3.3cm}@{}}
\toprule
\textbf{Thành phần} & \textbf{Vai trò} & \textbf{Quyết định nghiệp vụ} & \textbf{Credentials} \\
\midrule
\texttt{frontend} & \texttt{frontend} cho giáo viên và học sinh. Chỉ nói chuyện với \texttt{backend} & Không & Không có \\
\addlinespace
\texttt{backend} & Nơi giữ mọi quyết định nghiệp vụ, và là cửa vào duy nhất của hệ thống & \textbf{Mọi quyết định} & postgres / redis / minio \\
\addlinespace
\texttt{agent} & Kết nối với Model Provider. Thực hiện yêu cầu của giáo viên và giải đáp học sinh thông qua các cuộc trò chuyện & Không & Không có \\
\addlinespace
\texttt{document} & Đọc/Xử lí tài liệu giáo viên tải lên phục vụ \texttt{agent} sau này. & Không & minio (chỉ đọc) \\
\addlinespace
\texttt{ingest} & Đẩy dữ liệu từ \texttt(document) vào database --- nơi \texttt{agent} truy vấn thông qua \texttt{backend}  & Không & postgres (chỉ ghi) \\
\bottomrule
\end{tabular}
\end{table}

\texttt{frontend} không biết \texttt{agent} tồn tại. Mọi thứ nó cần đều đi qua \texttt{backend}, nên hệ thống
chỉ có \textbf{một cửa ra vào}.

\begin{landscape}
\begin{figure}[p]
  \centering
  \includegraphics[width=\linewidth,height=\textheight,keepaspectratio]{system-architecture.pdf}
  \caption[System Architecture]{System Architecture}
  \label{fig:system-architecture}
\end{figure}
\end{landscape}

\section{Frontend không nhìn thấy bên trong hệ thống}

Hình~\ref{fig:system-architecture} chia hệ thống làm hai vùng, và đặt \texttt{frontend} ở ngoài cả
hai. Trong VPC, \texttt{backend} nằm ở \emph{public subnet}; \texttt{agent}, \texttt{document},
\texttt{ingest} và ba kho nằm ở \emph{private subnet} --- chỉ những gì ở public subnet mới nhận
được kết nối từ bên ngoài.

\texttt{frontend} không bao giờ biết bên trong có mấy thành phần, hay chúng tên gì.

\section{Backend quyết định mọi thứ}

Đây là ràng buộc quan trọng nhất của kiến trúc, và nó là một ràng buộc \emph{nghiệp vụ} chứ không
phải một lựa chọn kỹ thuật. Luật chỉ có một câu: \textbf{\texttt{agent} nộp nội dung / đề xuất hành động,
\texttt{backend} quyết định làm gì.}

\texttt{agent} trả về \texttt{backend} hai thứ:

\begin{itemize}
  \item \textbf{Nội dung} --- câu hỏi và lời giải, một lượt trả lời cho giáo viên hay học
        sinh, một bản phân tích đọc từ dữ liệu mà \texttt{backend} cấp cho nó.
  \item \textbf{Đề xuất hành động} --- tool nào nên chạy, với tham số nào.
\end{itemize}

Hiệu lực do \texttt{backend} quyết định. Với nội dung: kiểm trước khi ghi.
Với hành động: chỉ chạy những tool chính nó cấp cho mô hình, và từ chối một đề xuất gọi tool
ngoài danh sách ấy hay thiếu tham số.

Chỗ thấy rõ nhất là việc soạn đề:

\begin{itemize}
  \item \texttt{agent} trả về một \textbf{đề xuất}: câu hỏi, các phương án, ánh xạ từ mỗi phương
        án nhiễu sang một lỗi cụ thể, và lời giải.
  \item \texttt{backend} \textbf{kiểm lại từng câu trước khi ghi} --- đúng một đáp án đúng, mọi phương
        án nhiễu đều gắn một lỗi, ít nhất hai cách giải --- rồi mới lưu. Và một đề chỉ rời trạng
        thái nháp khi giáo viên bấm duyệt.
\end{itemize}

\section{Thành phần nào kết nối tới kho nào?}

Hệ thống có ba kho --- kho dữ liệu chính (postgres), hàng đợi việc (redis), và kho tệp (minio).

\paragraph{Kho dữ liệu chính thuộc về \texttt{backend} và \texttt{ingest}.} Nó giữ lớp, học sinh, đề,
lần làm bài, sổ điểm ba mức, bộ đếm vòng, từng lượt làm lại kèm đề đã sinh ra, các đoạn hội
thoại, và các báo cáo. 

\paragraph{Kho tệp giữ \textbf{byte} của tài liệu, tức nội dung thô của tệp.} Đây là kho
duy nhất không phải cơ sở dữ liệu, và là kho duy nhất mà một thành phần chỉ đọc chứ không
bao giờ ghi: \texttt{backend} ghi vào, \texttt{document} đọc ra, và không có chiều ngược lại.

\texttt{agent} và \texttt{document} không có credential của kho dữ liệu chính.

\section{Các thành phần nói chuyện với nhau thế nào}
\label{sec:noi-chuyen}

Có \textbf{hai} cơ chế chở việc, và một tín hiệu rất ngắn không chở gì.

\paragraph{Một lời gọi, khi có người đang chờ.} \texttt{frontend} gọi \texttt{backend}. Đó là đường duy nhất
đi vào hệ thống từ bên ngoài, và cũng là lời gọi api duy nhất giữa hai thành phần.

\paragraph{Một việc trên hàng đợi, khi không ai chờ.} Mọi trao đổi còn lại đi đường này. Bên gửi
đặt một việc lên hàng đợi \texttt(redis) kèm \textbf{tên của việc} --- một chuỗi ký tự --- rồi đi tiếp. Bên nhận
lấy việc xuống và làm.

\paragraph{Và một tín hiệu rất ngắn, khi có trạng thái thay đổi.} Khi giáo viên tải lên một tài liệu, tài liệu hiển thị trạng thái
``đang xử lí'' trên \texttt{frontend}, lúc này \texttt{backend} đẩy
việc xử lí tài liệu lên hàng đợi \texttt{redis}, \texttt{document} làm xong thì đẩy tiếp một việc
nữa, và \texttt{ingest} lấy việc đó để ghi vào kho dữ liệu chính. Khi \texttt{ingest}
ghi xong một kết quả, nó đẩy một tín hiệu vào hàng đợi \texttt{redis}; \texttt{backend} nghe được thì đẩy xuống \texttt{frontend} hiển thị ``đã xong''.

\section{Việc gì đi qua hàng đợi, việc gì không}
\label{sec:hang-doi}

Mọi việc cần gọi mô hình đều chạy qua một hàng đợi vì một lần gọi mô hình lâu tới mức giữ một kết nối mở suốt thời gian chờ là không hợp lý.

\begin{table}[htbp]
\centering
\caption{Bốn việc thật sự cần gọi mô hình}
\label{tab:bon-viec}
\begin{tabular}{@{}L{4.2cm}L{9.6cm}@{}}
\toprule
\textbf{Việc} & \textbf{Khi nào xảy ra} \\
\midrule
Viết một câu hỏi & Mỗi câu của bộ đề giáo viên yêu cầu. \textbf{Một câu một việc}, không phải cả bộ một việc \\
\addlinespace
Sinh câu biến thể & Ngay khi học sinh nộp bài giai đoạn 1, để câu có sẵn lúc em bấm nút \\
\addlinespace
Trả lời một lượt hội thoại & Mỗi lượt học sinh hỏi lại ở giai đoạn 2 \\
\addlinespace
Đề xuất bước tiếp theo & Mỗi bước của một lượt trò chuyện với giáo viên \\
\bottomrule
\end{tabular}
\end{table}

\paragraph{Chấm bài không đi qua hàng đợi.} Với trắc nghiệm, chấm là một phép so giữa phương án
đã chọn và đáp án đúng đã lưu sẵn, nên nó chạy ngay trong lúc học sinh nộp bài. Kết quả nhìn thấy
được trên màn hình: giữa lúc nộp và lúc xem kết quả \textbf{không có trạng thái đang chấm nào}.

\paragraph{Màn hình được báo, thay vì phải hỏi lại.} Có những chỗ dữ liệu hiện ra dần:

\begin{itemize}
  \item \textbf{Hội thoại của học sinh} hiện chữ dần cho học sinh.
  \item \textbf{Hội thoại của giáo viên} đẩy từng bước của một lượt ngay khi bước đó xảy ra.
  \item \textbf{Thư viện tài liệu của giáo viên} lúc vừa tải lên hiển thị ``đang xử lí'', sau đó hiển thị ``đã xong'' mà không cần refresh màn hình.
\end{itemize}

\section{Nếu một thành phần ngừng thì mất gì?}

Một kiến trúc chia nhiều phần chỉ đáng giá nếu một phần hỏng không kéo cả hệ thống theo. Bảng
\ref{tab:suy-giam} nói rõ từng trường hợp.

\begin{table}[H]
\centering
\caption{Mỗi thành phần ngừng thì mất gì, còn gì}
\label{tab:suy-giam}
\begin{tabular}{@{}lL{5.0cm}L{5.6cm}@{}}
\toprule
\textbf{Ngừng} & \textbf{Mất} & \textbf{Còn} \\
\midrule
\texttt{redis} & Mọi việc cần gọi mô hình và tải lên tài liệu & Phát hành đề, làm bài, nộp bài, xem điểm giai đoạn 1 \\
\addlinespace
\texttt{minio} & \texttt{backend} không khởi động & --- \\
\addlinespace
\texttt{postgres} & \texttt{backend} không khởi động & --- \\
\addlinespace
\texttt{agent} & Soạn đề, hội thoại & Mọi thứ còn lại \\
\addlinespace
\texttt{document} & Tài liệu mới không tải lên được & Thư viện vẫn mở, tài liệu cũ vẫn dùng được \\
\addlinespace
\texttt{ingest} & Tài liệu mới không tải lên được & Việc \textbf{không mất}: nó nằm lại trên hàng đợi, bật lại thì chạy tiếp \\
\bottomrule
\end{tabular}
\end{table}

\section{Mô hình dữ liệu}

Mô hình dữ liệu có hai mươi mốt bảng, chia theo năm nhóm. Mỗi bảng trả lời đúng một câu hỏi
nghiệp vụ, và phần ghi chú dưới đây nói vì sao ranh giới giữa chúng nằm ở chỗ đó.

\begin{longtable}{@{}lL{9.1cm}@{}}
\caption{Hai mươi mốt bảng, theo năm nhóm}
\label{tab:bang-du-lieu}\\
\toprule
\textbf{Bảng} & \textbf{Giữ gì} \\
\midrule
\endfirsthead
\toprule
\textbf{Bảng} & \textbf{Giữ gì} \\
\midrule
\endhead
\midrule
\multicolumn{2}{r@{}}{\small\emph{tiếp trang sau}} \\
\endfoot
\bottomrule
\endlastfoot

\multicolumn{2}{@{}l@{}}{\textbf{Nhóm 1 --- người và lớp}} \\
\addlinespace
\texttt{classes} & Một lớp giáo viên tạo. Lớp thuộc về giáo viên, không thuộc về nhà trường. Tên lớp \emph{không} phải định danh \\
\cmidrule{1-2}
\texttt{students} & Họ tên, mã học sinh, thuộc lớp nào. Mã học sinh là khoá con người dùng \\
\cmidrule{1-2}
\texttt{teachers} & Họ tên và mã giáo viên \\
\addlinespace

\multicolumn{2}{@{}l@{}}{\textbf{Nhóm 2 --- nội dung đề}} \\
\addlinespace
\texttt{assessments} & Tác giả, tiêu đề, môn, khối, và \emph{trạng thái} --- bốn giá trị của vòng đời một đề: rỗng, có câu hỏi, đã duyệt, đã phát hành \\
\cmidrule{1-2}
\texttt{questions} & Đề bài, thứ tự câu học sinh nhìn thấy, mục tiêu học tập \\
\cmidrule{1-2}
\texttt{options} & Nhãn, nội dung, dấu đáp án đúng, và \emph{lỗi mà phương án này đại diện}. Ô lỗi để trống ở đúng một dòng mỗi câu, và dòng đó chính là đáp án đúng \\
\cmidrule{1-2}
\texttt{methods} & Các cách giải của một câu. Bắt buộc nhiều hơn một \\
\cmidrule{1-2}
\texttt{publications} & Sáu tham số phát hành, một dòng cho mỗi cặp đề và lớp, nên thu hồi được từng lớp một \\
\cmidrule{1-2}
\texttt{documents} & Tài liệu giáo viên tải lên: tên tệp, kích thước, khoá trỏ tới kho tệp, trạng thái xử lý, số trang, lý do hỏng. \textbf{Không giữ byte nào}. Thuộc về giáo viên, không thuộc về đề \\
\cmidrule{1-2}
\texttt{draft\_briefs} & Yêu cầu soạn của một đề: môn, khối, phạm vi, tổng số câu đã xin \\
\cmidrule{1-2}
\texttt{draft\_items} & Một ô cho mỗi câu đang soạn, kèm trạng thái, số lần đã thử, và \emph{lý do} lần gần nhất bị loại \\
\addlinespace

\multicolumn{2}{@{}l@{}}{\textbf{Nhóm 3 --- bài làm và điểm}} \\
\addlinespace
\texttt{attempts} & Một học sinh làm một đề: lớp lúc bắt đầu, giờ bắt đầu, giờ phải xong, giờ đã nộp \\
\cmidrule{1-2}
\texttt{answers} & Lựa chọn ở giai đoạn 1, ghi mỗi lần bấm chứ không gom tới lúc nộp \\
\cmidrule{1-2}
\texttt{question\_outcomes} & \textbf{Sổ điểm}: mức điểm cuối, lý do, số vòng đã dùng, đã chốt hay chưa \\
\addlinespace

\multicolumn{2}{@{}l@{}}{\textbf{Nhóm 4 --- giai đoạn 2}} \\
\addlinespace
\texttt{rounds} & Một lượt làm lại và giờ kết thúc của nó, đã cắt xuống hạn giai đoạn 2 nếu vượt \\
\cmidrule{1-2}
\texttt{round\_items} & Câu sinh cho lượt đó, kèm phương án và lời giải --- giữ cả đáp án đúng ở đây là thứ cho phép hệ thống tự chấm lượt \\
\cmidrule{1-2}
\texttt{pregenerated\_items} & Câu viết \emph{trước} khi học sinh bấm nút, trong lúc em còn đang đọc giải thích \\
\cmidrule{1-2}
\texttt{chat\_messages} & Từng lượt hội thoại, lưu \emph{trước} khi phát ra màn hình \\
\cmidrule{1-2}
\texttt{reports} & Báo cáo \emph{giải thích chưa rõ}, gắn với cả bài chứ không với một tin nhắn \\
\addlinespace

\multicolumn{2}{@{}l@{}}{\textbf{Nhóm 5 --- hội thoại của giáo viên}} \\
\addlinespace
\texttt{teacher\_conversations} & Một mạch hội thoại của một giáo viên. Tiêu đề do trợ lí đặt sau lượt đầu, giáo viên sửa được; xoá là \emph{xoá mềm} \\
\cmidrule{1-2}
\texttt{teacher\_turns} & Từng \emph{bước} của một lượt: lời nói, hoặc công cụ đã chạy, kèm chủ thể mà bước đó nói về \\
\end{longtable}
